% Options for packages loaded elsewhere
\PassOptionsToPackage{unicode}{hyperref}
\PassOptionsToPackage{hyphens}{url}
%
\documentclass[
]{article}
\usepackage{amsmath,amssymb}
\usepackage{lmodern}
\usepackage{iftex}
\ifPDFTeX
  \usepackage[T1]{fontenc}
  \usepackage[utf8]{inputenc}
  \usepackage{textcomp} % provide euro and other symbols
\else % if luatex or xetex
  \usepackage{unicode-math}
  \defaultfontfeatures{Scale=MatchLowercase}
  \defaultfontfeatures[\rmfamily]{Ligatures=TeX,Scale=1}
\fi
% Use upquote if available, for straight quotes in verbatim environments
\IfFileExists{upquote.sty}{\usepackage{upquote}}{}
\IfFileExists{microtype.sty}{% use microtype if available
  \usepackage[]{microtype}
  \UseMicrotypeSet[protrusion]{basicmath} % disable protrusion for tt fonts
}{}
\makeatletter
\@ifundefined{KOMAClassName}{% if non-KOMA class
  \IfFileExists{parskip.sty}{%
    \usepackage{parskip}
  }{% else
    \setlength{\parindent}{0pt}
    \setlength{\parskip}{6pt plus 2pt minus 1pt}}
}{% if KOMA class
  \KOMAoptions{parskip=half}}
\makeatother
\usepackage{xcolor}
\usepackage[margin=1in]{geometry}
\usepackage{color}
\usepackage{fancyvrb}
\newcommand{\VerbBar}{|}
\newcommand{\VERB}{\Verb[commandchars=\\\{\}]}
\DefineVerbatimEnvironment{Highlighting}{Verbatim}{commandchars=\\\{\}}
% Add ',fontsize=\small' for more characters per line
\usepackage{framed}
\definecolor{shadecolor}{RGB}{248,248,248}
\newenvironment{Shaded}{\begin{snugshade}}{\end{snugshade}}
\newcommand{\AlertTok}[1]{\textcolor[rgb]{0.94,0.16,0.16}{#1}}
\newcommand{\AnnotationTok}[1]{\textcolor[rgb]{0.56,0.35,0.01}{\textbf{\textit{#1}}}}
\newcommand{\AttributeTok}[1]{\textcolor[rgb]{0.77,0.63,0.00}{#1}}
\newcommand{\BaseNTok}[1]{\textcolor[rgb]{0.00,0.00,0.81}{#1}}
\newcommand{\BuiltInTok}[1]{#1}
\newcommand{\CharTok}[1]{\textcolor[rgb]{0.31,0.60,0.02}{#1}}
\newcommand{\CommentTok}[1]{\textcolor[rgb]{0.56,0.35,0.01}{\textit{#1}}}
\newcommand{\CommentVarTok}[1]{\textcolor[rgb]{0.56,0.35,0.01}{\textbf{\textit{#1}}}}
\newcommand{\ConstantTok}[1]{\textcolor[rgb]{0.00,0.00,0.00}{#1}}
\newcommand{\ControlFlowTok}[1]{\textcolor[rgb]{0.13,0.29,0.53}{\textbf{#1}}}
\newcommand{\DataTypeTok}[1]{\textcolor[rgb]{0.13,0.29,0.53}{#1}}
\newcommand{\DecValTok}[1]{\textcolor[rgb]{0.00,0.00,0.81}{#1}}
\newcommand{\DocumentationTok}[1]{\textcolor[rgb]{0.56,0.35,0.01}{\textbf{\textit{#1}}}}
\newcommand{\ErrorTok}[1]{\textcolor[rgb]{0.64,0.00,0.00}{\textbf{#1}}}
\newcommand{\ExtensionTok}[1]{#1}
\newcommand{\FloatTok}[1]{\textcolor[rgb]{0.00,0.00,0.81}{#1}}
\newcommand{\FunctionTok}[1]{\textcolor[rgb]{0.00,0.00,0.00}{#1}}
\newcommand{\ImportTok}[1]{#1}
\newcommand{\InformationTok}[1]{\textcolor[rgb]{0.56,0.35,0.01}{\textbf{\textit{#1}}}}
\newcommand{\KeywordTok}[1]{\textcolor[rgb]{0.13,0.29,0.53}{\textbf{#1}}}
\newcommand{\NormalTok}[1]{#1}
\newcommand{\OperatorTok}[1]{\textcolor[rgb]{0.81,0.36,0.00}{\textbf{#1}}}
\newcommand{\OtherTok}[1]{\textcolor[rgb]{0.56,0.35,0.01}{#1}}
\newcommand{\PreprocessorTok}[1]{\textcolor[rgb]{0.56,0.35,0.01}{\textit{#1}}}
\newcommand{\RegionMarkerTok}[1]{#1}
\newcommand{\SpecialCharTok}[1]{\textcolor[rgb]{0.00,0.00,0.00}{#1}}
\newcommand{\SpecialStringTok}[1]{\textcolor[rgb]{0.31,0.60,0.02}{#1}}
\newcommand{\StringTok}[1]{\textcolor[rgb]{0.31,0.60,0.02}{#1}}
\newcommand{\VariableTok}[1]{\textcolor[rgb]{0.00,0.00,0.00}{#1}}
\newcommand{\VerbatimStringTok}[1]{\textcolor[rgb]{0.31,0.60,0.02}{#1}}
\newcommand{\WarningTok}[1]{\textcolor[rgb]{0.56,0.35,0.01}{\textbf{\textit{#1}}}}
\usepackage{graphicx}
\makeatletter
\def\maxwidth{\ifdim\Gin@nat@width>\linewidth\linewidth\else\Gin@nat@width\fi}
\def\maxheight{\ifdim\Gin@nat@height>\textheight\textheight\else\Gin@nat@height\fi}
\makeatother
% Scale images if necessary, so that they will not overflow the page
% margins by default, and it is still possible to overwrite the defaults
% using explicit options in \includegraphics[width, height, ...]{}
\setkeys{Gin}{width=\maxwidth,height=\maxheight,keepaspectratio}
% Set default figure placement to htbp
\makeatletter
\def\fps@figure{htbp}
\makeatother
\setlength{\emergencystretch}{3em} % prevent overfull lines
\providecommand{\tightlist}{%
  \setlength{\itemsep}{0pt}\setlength{\parskip}{0pt}}
\setcounter{secnumdepth}{-\maxdimen} % remove section numbering
\usepackage{statstyle}
\ifLuaTeX
  \usepackage{selnolig}  % disable illegal ligatures
\fi
\IfFileExists{bookmark.sty}{\usepackage{bookmark}}{\usepackage{hyperref}}
\IfFileExists{xurl.sty}{\usepackage{xurl}}{} % add URL line breaks if available
\urlstyle{same} % disable monospaced font for URLs
\hypersetup{
  pdftitle={Simulation study for debiased learning},
  hidelinks,
  pdfcreator={LaTeX via pandoc}}

\title{Simulation study for debiased learning}
\author{}
\date{\vspace{-2.5em}}

\begin{document}
\maketitle

Let \(A\) be a sample of size \(n =\) 500 selected from the population
\(U\) of size \(N =\) 5000. Let \(y\) be a variable of interest and we
try to estimate the population total \(t_y = \sum_{i \in U} y_i\) and
\(\bf x_{i}\) of size \(p\) be the auxiliary variables associated with
\(y\) where \(x_{ij} \iid N(2, 1)\).
\textcolor{red}{Make sure $X$ are away from zero. So, we may use $X \sim N(3,1)$ }

Suppose the following
superpopulation model holds \[
y_i = \bf x_i^T \bf \beta + \varepsilon_i, \quad \varepsilon_i \iid N(0,\sigma^2), i = 1, \cdots, N
\] where \(\sigma^2 = 1\) and
\(\bf \beta = (1, \cdots, 1, 0, \cdots, 0)\) with \(s =\) 5 1s and
\(p - s =\) 195 0s. We assume the samples are drawn from the \(\pi\)-ps
sampling with first order probability proportional to \[
\pi_i \propto x_{i1}^2
\]

\textcolor{red}{We can use SRS for initial simulation study. Considering unequal sampling will only complicate the setup and is not really helpful in understanding the problem. 
}

We consider the following estimators.

\begin{itemize}
  \item Difference estimator
  $$
  \hat t_y = \pa{\sum_{i \in U}\bf x_i}^T \bf \beta + \sum_{i \in A} d_i \pa{y_i - \bf x_i^T \bf \beta}
  $$
  \item GREG estimator
  $$
  \hat t_y = \pa{\sum_{i \in U}\bf x_i}^T \hat{\bf \beta} + \sum_{i \in A} d_i \pa{y_i - \bf x_i^T \hat{\bf \beta}}
  $$
  where $\hat {\bf \beta} = \pa{\sum_{i \in A} d_i \bf x_i \bf x_i^T}^{-1} \pa{\sum_{i \in A} d_i \bf x_i y_i}$
  
  \item Model-based estimator
  $$
  \hat t_y = \pa{\sum_{i \in U}\bf x_i}^T \hat{\bf \beta}
  $$
  
  \item Debiased GREG estimator(Proposed)
    $$
  \hat t_y = \pa{\sum_{i \in U}\bf x_i}^T \pa{ \frac{\sum_{i \in A_2}d_i}{\sum_{i \in A}d_i} \hat{\bf \beta}_1 + \frac{\sum_{i \in A_1}d_i}{\sum_{i \in A}d_i} \hat{\bf \beta}_2} + \sum_{i \in A_1} d_i \pa{y_i - \bf x_i^T \hat{\bf \beta}_2} + \sum_{i \in A_2} d_i \pa{y_i - \bf x_i^T \hat{\bf \beta}_1}
  $$
  where $A_1$ and $A_2$ are the subsamples of $A$ with $A_1 \bigcup A_2 = A$, and $\hat {\bf \beta}_1$ and $\hat {\bf \beta}_2$ are the regression coefficients built upon the sample $A_1$ and sample $A_2$ respectively.
  \item Design unbiased GREG estimator(Sande, 2021)
  $$
  \hat t_y = \pa{\sum_{i \in U}\bf x_i}^T \hat{\bf \beta}_1 + \sum_{i \in A_1} \pa{y_i - \bf x_i^T \hat{\bf \beta}_1} + \sum_{i \in A_2} \frac{1}{P(i \in A_2 | i \in A)} \pa{y_i - \bf x_i^T \hat{\bf \beta}_1}
  $$
  
  %\item LASSO estimator
  %\item Debiased Lasso estimator(Zhang \& Zhang, 2014)
\end{itemize}

\begin{Shaded}
\begin{Highlighting}[]
\NormalTok{res }\OtherTok{=} \ConstantTok{NULL}
\ControlFlowTok{for}\NormalTok{(simnum }\ControlFlowTok{in} \DecValTok{1}\SpecialCharTok{:}\DecValTok{500}\NormalTok{)\{}
  \FunctionTok{set.seed}\NormalTok{(simnum)}
\NormalTok{  Index }\OtherTok{=} \FunctionTok{sample}\NormalTok{(}\DecValTok{1}\SpecialCharTok{:}\NormalTok{N, }\AttributeTok{size =}\NormalTok{ n, }\AttributeTok{replace =} \ConstantTok{FALSE}\NormalTok{, }\AttributeTok{prob =}\NormalTok{ pi1)}
\NormalTok{y\_s }\OtherTok{=}\NormalTok{ y[Index]}
\NormalTok{X\_s }\OtherTok{=}\NormalTok{ X[Index, ]}
\NormalTok{d1\_s }\OtherTok{=}\NormalTok{ d1[Index]}
\NormalTok{lm\_obj }\OtherTok{=} \FunctionTok{lm}\NormalTok{(y\_s }\SpecialCharTok{\textasciitilde{}}\NormalTok{  X\_s, }\AttributeTok{weights =}\NormalTok{ d1\_s)}
\NormalTok{beta\_hat }\OtherTok{=}\NormalTok{ lm\_obj}\SpecialCharTok{$}\NormalTok{coefficients[}\SpecialCharTok{{-}}\DecValTok{1}\NormalTok{]}
\CommentTok{\#drop(solve(t(X\_s) \%*\% diag(d1\_s) \%*\% X\_s, t(X\_s) \%*\% diag(d1\_s) \%*\% y\_s))}

\CommentTok{\# beta\_hat = unname(lm(y\_s \textasciitilde{}  t(apply(X\_s, 1, function(k) k }
\CommentTok{\# {-} colSums(X\_s * d1\_s) / N)), weights = d1\_s)$coefficients[{-}1])}

\NormalTok{y\_diff }\OtherTok{=} \FunctionTok{drop}\NormalTok{(}\FunctionTok{colSums}\NormalTok{(X) }\SpecialCharTok{\%*\%}\NormalTok{ beta }\SpecialCharTok{+} \FunctionTok{sum}\NormalTok{((y\_s }\SpecialCharTok{{-}} \FunctionTok{drop}\NormalTok{(X\_s }\SpecialCharTok{\%*\%}\NormalTok{ beta)) }\SpecialCharTok{*}\NormalTok{ d1\_s))}

\NormalTok{y\_GREG }\OtherTok{=} \FunctionTok{drop}\NormalTok{(}\FunctionTok{colSums}\NormalTok{(X) }\SpecialCharTok{\%*\%}\NormalTok{ beta\_hat }\SpecialCharTok{+} \FunctionTok{sum}\NormalTok{((y\_s }\SpecialCharTok{{-}} \FunctionTok{drop}\NormalTok{(X\_s }\SpecialCharTok{\%*\%}\NormalTok{ beta\_hat)) }\SpecialCharTok{*}\NormalTok{ d1\_s))}

\NormalTok{y\_model }\OtherTok{=} \FunctionTok{drop}\NormalTok{(}\FunctionTok{colSums}\NormalTok{(X) }\SpecialCharTok{\%*\%}\NormalTok{ beta\_hat)}

\NormalTok{Index\_sub1 }\OtherTok{=}\NormalTok{ Index[}\FunctionTok{sample}\NormalTok{(}\DecValTok{1}\SpecialCharTok{:}\NormalTok{n, }\AttributeTok{size =}\NormalTok{ m, }\AttributeTok{replace =} \ConstantTok{FALSE}\NormalTok{)]}
\NormalTok{Index\_sub2 }\OtherTok{=}\NormalTok{ Index[}\SpecialCharTok{!}\NormalTok{(Index }\SpecialCharTok{\%in\%}\NormalTok{ Index\_sub1)]}

\NormalTok{y\_s1 }\OtherTok{=}\NormalTok{ y[Index\_sub1]}
\NormalTok{X\_s1 }\OtherTok{=}\NormalTok{ X[Index\_sub1, ]}
\NormalTok{X1\_s1 }\OtherTok{=} \FunctionTok{cbind}\NormalTok{(}\DecValTok{1}\NormalTok{, X)[Index\_sub1, ]}
\NormalTok{d1\_s1 }\OtherTok{=}\NormalTok{ d1[Index\_sub1]}
\NormalTok{lm\_obj1 }\OtherTok{=} \FunctionTok{lm}\NormalTok{(y\_s1 }\SpecialCharTok{\textasciitilde{}}\NormalTok{  X\_s1, }\AttributeTok{weights =}\NormalTok{ d1\_s1)}
\NormalTok{beta\_hat1 }\OtherTok{=}\NormalTok{ lm\_obj1}\SpecialCharTok{$}\NormalTok{coefficients[}\SpecialCharTok{{-}}\DecValTok{1}\NormalTok{]}

\NormalTok{y\_s2 }\OtherTok{=}\NormalTok{ y[Index\_sub2]}
\NormalTok{X\_s2 }\OtherTok{=}\NormalTok{ X[Index\_sub2, ]}
\NormalTok{d1\_s2 }\OtherTok{=}\NormalTok{ d1[Index\_sub2]}
\NormalTok{lm\_obj2 }\OtherTok{=} \FunctionTok{lm}\NormalTok{(y\_s2 }\SpecialCharTok{\textasciitilde{}}\NormalTok{  X\_s2, }\AttributeTok{weights =}\NormalTok{ d1\_s2)}
\NormalTok{beta\_hat2 }\OtherTok{=}\NormalTok{ lm\_obj2}\SpecialCharTok{$}\NormalTok{coefficients[}\SpecialCharTok{{-}}\DecValTok{1}\NormalTok{]}

\CommentTok{\# y\_debiased = drop(colSums(X) \%*\% beta\_hat + sum((y\_s1 {-} drop(X\_s1 \%*\% beta\_hat2)) * d1\_s1) }
\CommentTok{\# + sum((y\_s2 {-} drop(X\_s2 \%*\% beta\_hat1)) * d1\_s2))}

\NormalTok{y\_debiased }\OtherTok{=} \FunctionTok{drop}\NormalTok{(}\FunctionTok{colSums}\NormalTok{(X) }\SpecialCharTok{\%*\%}\NormalTok{ (}\FunctionTok{sum}\NormalTok{(d1\_s2) }\SpecialCharTok{/} \FunctionTok{sum}\NormalTok{(d1\_s) }\SpecialCharTok{*}\NormalTok{ beta\_hat1 }\SpecialCharTok{+} \FunctionTok{sum}\NormalTok{(d1\_s1) }\SpecialCharTok{/} \FunctionTok{sum}\NormalTok{(d1\_s) }\SpecialCharTok{*}\NormalTok{ beta\_hat2) }\SpecialCharTok{+} 
                    \FunctionTok{sum}\NormalTok{((y\_s1 }\SpecialCharTok{{-}} \FunctionTok{drop}\NormalTok{(X\_s1 }\SpecialCharTok{\%*\%}\NormalTok{ beta\_hat2)) }\SpecialCharTok{*}\NormalTok{ d1\_s1) }\SpecialCharTok{+} \FunctionTok{sum}\NormalTok{((y\_s2 }\SpecialCharTok{{-}} \FunctionTok{drop}\NormalTok{(X\_s2 }\SpecialCharTok{\%*\%}\NormalTok{ beta\_hat1)) }\SpecialCharTok{*}\NormalTok{ d1\_s2))}

\NormalTok{y\_unbiased }\OtherTok{=} \FunctionTok{drop}\NormalTok{(}\FunctionTok{colSums}\NormalTok{(X) }\SpecialCharTok{\%*\%}\NormalTok{ beta\_hat1 }\SpecialCharTok{+} \FunctionTok{sum}\NormalTok{((y\_s1 }\SpecialCharTok{{-}} \FunctionTok{drop}\NormalTok{(X\_s1 }\SpecialCharTok{\%*\%}\NormalTok{ beta\_hat1))) }
                  \SpecialCharTok{+} \FunctionTok{sum}\NormalTok{((y\_s2 }\SpecialCharTok{{-}} \FunctionTok{drop}\NormalTok{(X\_s2 }\SpecialCharTok{\%*\%}\NormalTok{ beta\_hat1)) }\SpecialCharTok{*} \DecValTok{2}\NormalTok{))}

\NormalTok{res }\OtherTok{=} \FunctionTok{rbind}\NormalTok{(res, }\FunctionTok{c}\NormalTok{(}\AttributeTok{Diff =}\NormalTok{ y\_diff, }\AttributeTok{GREG =}\NormalTok{ y\_GREG, }\AttributeTok{Model =}\NormalTok{ y\_model, }\AttributeTok{Debiased =}\NormalTok{ y\_debiased, }\AttributeTok{Unbiased =}\NormalTok{ y\_unbiased))}
\NormalTok{\}}

\NormalTok{BIAS }\OtherTok{=} \FunctionTok{colMeans}\NormalTok{(res }\SpecialCharTok{{-}}\NormalTok{ t\_y)}
\NormalTok{SE }\OtherTok{=} \FunctionTok{apply}\NormalTok{(res, }\DecValTok{2}\NormalTok{, }\ControlFlowTok{function}\NormalTok{(x) }\FunctionTok{sqrt}\NormalTok{(}\FunctionTok{var}\NormalTok{(x) }\SpecialCharTok{*}\NormalTok{ (}\FunctionTok{length}\NormalTok{(x)}\SpecialCharTok{{-}}\DecValTok{1}\NormalTok{)}\SpecialCharTok{/}\FunctionTok{length}\NormalTok{(x) ))}
\NormalTok{RMSE }\OtherTok{=} \FunctionTok{apply}\NormalTok{(res }\SpecialCharTok{{-}}\NormalTok{ t\_y, }\DecValTok{2}\NormalTok{, }\ControlFlowTok{function}\NormalTok{(x) }\FunctionTok{sqrt}\NormalTok{(}\FunctionTok{mean}\NormalTok{(x}\SpecialCharTok{\^{}}\DecValTok{2}\NormalTok{)))}

\FunctionTok{cbind}\NormalTok{(BIAS, SE, RMSE)}
\end{Highlighting}
\end{Shaded}

\begin{verbatim}
##                BIAS         SE       RMSE
## Diff      -23.73050   741.4037   741.7834
## GREG      117.23454  1263.6877  1269.1141
## Model    -900.42621  9127.5889  9171.8944
## Debiased   60.32471  2837.3919  2838.0331
## Unbiased -211.53455 16899.2166 16900.5405
\end{verbatim}

{\bf Additional comments about the simulation study}
\begin{itemize}
\item Using the high dimensional regression is OK for now, but we are more interested in nonparametric regression problem. We may consider using smoothing spline or Kernel ridge regression or deep learning method. 
\item The proposed method will be less biased but will have a higher variance. Variance estimation (and inference) should be another big advantage of the proposed method. 
\end{itemize}

\end{document}
